% Generated draft: requires mathematical and source review.
\documentclass[11pt]{article}
\usepackage[margin=25mm]{geometry}
\usepackage{fontspec}
\setmainfont{Latin Modern Roman}
\usepackage{amsmath,amssymb,mathtools,unicode-math}
\setmathfont{Latin Modern Math}
\usepackage{enumitem,needspace,xurl,hyperref}
\setlength{\parindent}{0pt}
\setlength{\parskip}{5pt}
\setlength{\emergencystretch}{3em}
\newcommand{\N}{\mathbb{N}}
\newcommand{\Z}{\mathbb{Z}}
\newcommand{\Q}{\mathbb{Q}}
\newcommand{\R}{\mathbb{R}}
\newcommand{\C}{\mathbb{C}}
\newcommand{\F}{\mathbb{F}}
\newcommand{\A}{\mathbb{A}}
\newcommand{\PP}{\mathbb{P}}
\newcommand{\E}{\mathbb{E}}
\newcommand{\eps}{\varepsilon}
\newcommand{\dd}{\,\mathrm{d}}
\begin{document}
\section*{UnsolvedMath ID: 30006560}
\textbf{Source problem code:} catalogue entry
\subsection{Definitions and mathematical statement}
Let \(r\in\Z\) with \(r\geq3\), and let \(R,G,B\in\Z_{\geq0}\). Under the usual finite simple hypergraph interpretation, let \(H=(V,E)\), where \(V\) is finite and \(E\subseteq\binom{V}{r-1}\). A three-coloring is a map \(c:E\to\{\mathrm{red},\mathrm{green},\mathrm{blue}\}\), so each edge receives exactly one color. Write \(E_{\mathrm{red}},E_{\mathrm{green}},E_{\mathrm{blue}}\) for the three color classes, and require
\[
|E_{\mathrm{red}}|=R,\qquad |E_{\mathrm{green}}|=G,\qquad |E_{\mathrm{blue}}|=B.
\]
For \(S\subseteq V\), the induced colored hypergraph \(H[S]\) has edge set \(\{e\in E:e\subseteq S\}\), with colors inherited from \(c\). Define
\[
N_r(H,c)=\left|\left\{S\in\binom{V}{r}:\ \text{for every }a\in\{\mathrm{red},\mathrm{green},\mathrm{blue}\}\text{ there exists }e\in E_a\text{ with }e\subseteq S\right\}\right|.
\]
The problem is to determine, for every such \(r,R,G,B\), the extremal quantity
\[
M_r(R,G,B)=\max_{(H,c)}N_r(H,c),
\]
where the maximum ranges over colored hypergraphs satisfying the preceding requirements, with no prescribed number of vertices. The finite simple hypergraph convention is an explicit interpretation: the supplied statement does not separately specify finiteness, simplicity, or whether zero-sized color classes are permitted.

The original excerpt additionally asserts that the answer is monotonically non-decreasing in \(r\); with the edge counts held fixed, this reads \(M_{r+1}(R,G,B)\geq M_r(R,G,B)\). This is recorded as a supplied contextual assertion, not independently verified here.

The supplied literature assessment reports an upper bound \(6\sqrt{RGB}\) for the count and a sharp upper bound \(2\sqrt{RGB}\) when \(r=3\). It also describes a conjectured bound for general \(r\) as unresolved, but does not explicitly state that conjecture. These are supplied dated literature claims, not independently verified results of this draft. In particular, the precise general conjecture cannot be recovered unambiguously from the assessment alone.
\subsection{Short English statement}
For each integer \(r\geq3\) and prescribed edge counts \(R,G,B\), determine the largest possible number of \(r\)-element vertex sets containing a red edge, a green edge, and a blue edge of a three-colored \((r-1)\)-uniform hypergraph. The formulation uses an explicitly stated finite simple hypergraph convention. The supplied record labels the problem open; its literature claims and verification labels require independent review.
\subsection{Sources}
\begin{itemize}
\item Supplied local UnsolvedMath registry record, permanent ID \(30006560\), title ``Vertex Sets Meeting Every Edge Color,'' source problem code \texttt{catalogue entry}, in the Oberwolfach Reports Open Problems collection. The record gives proposed year 2026 and status \texttt{open}.
\item Supplied report citation: \emph{Combinatorics (2026)}, Oberwolfach Reports, DOI 10.4171/owr/2026/1, pp.\ 5--79. Source URL: \url{https://doi.org/10.4171/owr/2026/1}. The supplied background describes Oberwolfach Reports as documenting mathematical workshops and their open-problem discussions.
\item Supplied literature-source entry: ``Oberwolfach Report 1/2026,'' \url{https://doi.org/10.4171/owr/2026/1}. The exact supplied literature-check date is \texttt{2026-08-22}. Its open-status assessment, upper-bound claims, sharpness claim for \(r=3\), and unresolved-conjecture assessment are supplied dated curation, not independently verified here. The background expressly requests independent verification of time-sensitive status and cited claims.
\item Attribution qualification: the original excerpt begins with the name Hans Yu and says that the question is motivated by incidence geometry. The structured \texttt{proposed\_by} field is null, so the precise attribution role of that name remains unresolved.
\item Source-verification qualification: the registry labels the statement \texttt{corrected\_verified} and reports that the official report was checked, the question isolated from surrounding exposition, and notation repaired. These are supplied registry assertions; this draft does not claim to have visited or checked the report.
\end{itemize}
\end{document}
